% !TEX root = chapter-standalone.tex

\section{Motivations}
\label{sec:simplest-dyn-system}

\todotext{Expand this intro a bit to explain the conceptual idea of actions better, and mention the key term ``state space''. Maybe move the robot on a grid example up to this intro as a concrete motivating example. Also, maybe we can motivate/explain why it makes sense to abstractly separate the set of ``instructions'' from the functions that model the dynamics... why is this a good idea? what are the benefits? maybe a second example can also underscore this, maybe one from linear algebra?}

\todotext{Because of our semicolon infix notation for function composition, the `more natural' version of actions to prefer is right actions rather than left actions (normally it is the other way around). So maybe we should start with right actions and prefer/emphasize them in our exposition (while of course still mentioning left actions and explaining them...)}

\todotext{This section is to be rewritten, with the 'covariant' (could we also say coordinate free?) version of the definition of action as the main one. Then what is currently called a right action should be 'braided' and then called a left action (but as a remark, no longer as a definition).}

\todotext{Add a remark that explains the term/concept of 'contravariant' action (there is material on this in devel at the moment. If we use the term 'contravarian morphism' somewhere, be sure we define it.)}

\todotext{Throughout this chapter, maybe we should use S as a default letter for the 'state set' rather A?}



The idea of ``actions'' appears across mathematics and in many areas of application -- for example in the topic of dynamical systems in engineering, or in the topic of automata in computer science. 

Here is the central pattern captured by the concept of an action: on the one hand, a monoid (or some other algebraic gadget) is given which models time or models a sequence of events, and on the other hand that monoid ``acts'' on a set that models the possible states of a system. This action encodes how the system changes over time or changes in reaction to a sequence of events or inputs. Before fitting this idea into mathematical definitions, let us begin with two motivating examples. 

   
\subsection{Motivating example: a robot on a grid}
   
    Consider a robot moving on a two-dimensional grid made up of 16 squares, as illustrated in \cref{fig:robot_example}. To label each of the 16 squares we use simple coordinates: labels that consist of an ``$x_1$-direction'' (horizontal) and an ``$x_2$-direction" (vertical), with each dimension having possible entries $1$, $2$, $3$ or $4$. In other words, the set of 16 squares is labeled by the set
\begin{equation}
    \setX = \setX_1 \cartprod \setX_2
\end{equation}
with $\setX_1 = \makeset{1, 2, 3, 4}$ and $\setX_2 = \makeset{1, 2, 3, 4}$. 

    \begin{marginfigure}
        \centering
        \includesag{120_robotcategory}
        \caption{Robot on a grid example. }
        \label{fig:robot_example}
    \end{marginfigure}

    We consider the position coordinates of the robot to be its ``state''.  For example, in the figure, the state of the robot is that its position is $\tup{1,1} \setin \setX$. 
    
    
    Next, consider the following set of instructions that we may give the robot,
\begin{equation}
\prin=\{\uparrow, \downarrow,\rightarrow,\leftarrow\},
\end{equation}
where each symbol is a command to ``move one square over'' in the corresponding direction  -- unless the movement would lead the robot to leave the grid, in which case no movement is undertaken. 


We can model this situation using the following  ``dynamics function'': 
    \begin{equation}
        \begin{aligned}
            \prdyn\colon \prin \cartprod \prst                & \sto \prst \\
            \tupp{\uparrow, \tup{\prsteln{1},\prsteln{2}}}    & \mapsto \tupp{\prsteln{1}, \min(\prsteln{2}+1,4)}, \\
            \tupp{\downarrow, \tup{\prsteln{1},\prsteln{2}}}  & \mapsto \tupp{\prsteln{1}, \max(\prsteln{2}-1,1)}, \\
            \tupp{\rightarrow, \tup{\prsteln{1},\prsteln{2}}} & \mapsto \tupp{\min(\prsteln{1}+1,4), \prsteln{2}}, \\
            \tupp{\leftarrow, \tup{\prsteln{1},\prsteln{2}}}  & \mapsto \tupp{\max(\prsteln{1}-1,1), \prsteln{2}}, \\
        \end{aligned}
    \end{equation}
    where the ``$\min$'' and ``$\max$'' ensure that the robot stays in within the boundaries of the grid.
    
    \todotext{Rework the above example so that it fits better with the action definition as we now have it...}
    

\devel{
\subsection{Motivating example: a discrete-time dynamical system}

\todotext{What is a nice motivating example here??}
}


\devel{
\subsection{Scaling of vectors}

\begin{equation}
    \defmapunnamed{
        \reals \cartprod \reals^n
    }{
        \to
    }{
        \reals^n
    }{
        \tup{\lambda, \mat{v}}
    }{
        \lambda\mat{v}
    }
\end{equation}

\begin{equation}
    (\lambda \mu)\mat{v} =  \lambda (\mu \mat{v})
\end{equation}

\begin{equation}
    1 \mat{v} = \mat{v}
\end{equation}

\subsection{Matrix-vector multiplication}

\begin{equation}
    \defmapunnamed{
        M_n(\reals) \cartprod \reals^n
    }{
        \to
    }{
        \reals^n
    }{
        \tup{\mat{A}, \mat{v}}
    }{
        \mat{A}\mat{v}
    }
\end{equation}

\begin{equation}
    (\mat{A} \mat{B})\mat{v} =  \mat{A} (\mat{B} \mat{v})
\end{equation}

\begin{equation}
    \mat{1}\mat{v} = \mat{v}
\end{equation}

\begin{equation}
    \defmapunnamed{
        GL_n(\reals) \cartprod \reals^n
    }{
        \to
    }{
        \reals^n
    }{
        \tup{\mat{A}, \mat{v}}
    }{
        \mat{A}\mat{v}
    }
\end{equation}

\subsection{Matrix-matrix multiplication}

\begin{equation}
    \defmapunnamed{
        M_n(\reals) \cartprod M_n(\reals)
    }{
        \to
    }{
        M_n(\reals)
    }{
        \tup{\mat{A}, \mat{B}}
    }{
        \mat{A}\mat{B}
    }
\end{equation}

\begin{equation}
    (\mat{A} \mat{B})\mat{C} =  \mat{A} (\mat{B} \mat{C})
\end{equation}

\begin{equation}
    \mat{1}\mat{A} = \mat{A}
\end{equation}

\subsection{Monoid multiplication}

\begin{equation}
    \defmapunnamed{
        \monoidAset \cartprod \monoidAset
    }{
        \to
    }{
        \monoidAset
    }{
        \tup{\monela, \monelb}
    }{
        \monela \mthen_\monoidA \monelb
    }
\end{equation}

\begin{equation}
    (\monela \mthen_\monoidA \monelb) \mthen_\monoidA \monelc =   \monela \mthen_\monoidA (\monelb \mthen_\monoidA \monelc)
\end{equation}

\begin{equation}
    \idmon_\monoidA \mthen_\monoidA \monela =  \monela
\end{equation}

\subsection{Groups of symmetries}

\begin{equation}
    \defmapunnamed{
        \grpAset \cartprod \setA
    }{
        \to
    }{
        \setA
    }{
        \tup{g, \ela}
    }{
        g \cdot \ela
    }
\end{equation}

\subsection{Flow maps}

\begin{equation}
    \defmapperiod{
        \varphi
    }{
        \reals \cartprod \reals^n
    }{
        \to
    }{
        \reals^n
    }{
        \tup{t, x_0}
    }{
        \varphi(t, x_0)
    }
\end{equation}

\subsection{Discrete dynamical systems}

\begin{equation}
    \defmapperiod{
        \varphi
    }{
        \wnumbers \cartprod X
    }{
        \to
    }{
        X
    }{
        \tup{n, x_0}
    }{
        \varphi(n, x_0)
    }
\end{equation}

\begin{equation}
    \defmapperiod{
        \act
    }{
        \setA \cartprod \wnumbers
    }{
        \to
    }{
        \setA
    }{
        \tup{\ela, n}
    }{
        \mapa^n(\ela)
    }
\end{equation}

}

